\subsection{局部凸初始拓扑}

命题 4. --- 设 E 是向量空间，\((\mathrm{E}_{\mathrm{t}})_{\mathrm{t}\in\mathrm{I}}\) 是一族局部凸空间。对每个 \(\iota\in I\) ，设 \(f_{t}\) 是 E 到 \(E_{t}\) 中的线性映射；则 E 上使每个映射 \(f_{\iota}\) 都连续的最粗拓扑 T 是一个局部凸拓扑。

利用 II, 第 24 页, 推论，这是半范数族定义的拓扑的相应性质 (II, 第 5 页) 的一个特殊情形。

特别地，局部凸空间的每个向量子空间，以及局部凸空间的每个乘积空间，都是局部凸的。局部凸空间的每个投射极限是局部凸的。

Fréchet 空间的每个可数乘积（特别地，Banach 空间的每个可数乘积）是 Fréchet 空间。

每个局部凸 Hausdorff 空间 E 同构于某个 Banach 空间乘积的一个子空间，且当 E 完备时该子空间是闭的 (II, 第 5 页, prop. 3)。每个 Fréchet 空间同构于 Banach 空间的一个可数乘积的闭子空间 (loc. cit.)。

